\section{Introduction}
\label{sec:intro}

% Paragraph 1: Motivate the stakes of blaming and defending as two directions of dissent.
When online communities judge everyday interpersonal conflicts, the majority's verdict is not the only perspective at stake. A dissenting judgment can call for accountability where others see no wrongdoing, or question condemnation that others accept. How these judgments are received matters for understanding whether a discussion accommodates challenges to both collective blame and collective defense. Neither kind of dissent is necessarily correct, but treating them as interchangeable obscures which challenges attract approval, opposition, or sustained conversation. We therefore focus on minority judgments about whether someone is at fault, distinguishing \emph{blaming dissent}, which assigns fault to someone whom most others defend, from \emph{defending dissent}, which rejects fault assigned by most others. Does the direction of dissent shape how others respond and how minority judgments develop over time?

% Paragraph 2: Establish alignment and agreement strength as the context for dissent.
Understanding this distinction requires situating each judgment within the surrounding distribution of opinion. A blaming judgment can express the majority position in one discussion and constitute a minority position in another. The strength of agreement also matters: opposing a narrow majority in a divided discussion differs from challenging a judgment that nearly everyone else supports. We therefore distinguish \emph{majority} from \emph{minority} judgments, and \emph{high-agreement} from \emph{low-agreement} posts. Our primary focus is minority dissent in high-agreement posts, where one judgment strongly predominates. These discussions make the question particularly salient: an overwhelming majority can coexist with continuing disagreement, and the final verdict alone tells us little about how that disagreement is treated or expressed. Low-agreement posts provide a comparison condition for interpreting this experience, rather than an equivalent setting of dissent against a strong majority.

% Paragraph 3: Use prior work and the spiral of silence to motivate the research gap.
Prior research provides important foundations for studying this experience. Work on Reddit shows that negative evaluations can coexist with attention through replies \cite{davis2021emotional}, while research on moral-judgment discussions examines how the value placed on agreement and dissent varies with consensus strength \cite{mellody2025consensus}. We revisit majority--minority differences in scores and conversational engagement across agreement levels as an empirical baseline, rather than treating their reappearance as a new research question. Our central question is what this aggregate comparison can conceal: whether the reception of dissent differs according to whether it blames or defends the person under judgment.

The spiral of silence further motivates examining minority participation alongside its reception. This account links perceptions of an unfavorable opinion climate to reluctance to speak out, through concern about social isolation \cite{noelleneumann1974spiral}. It directs attention to whether minority judgments continue to appear as a discussion unfolds, without implying that all dissent must steadily disappear. Our observational data capture public comments, not unspoken opinions, perceived isolation, or individual willingness to speak. We therefore examine changes in the prevalence of expressed minority judgments, rather than treating them as direct measurements of silence or self-censorship.

% Paragraph 4: Introduce AITA, the baseline, and the two primary research questions.
We investigate these questions primarily in r/AmItheAsshole (AITA), where an original poster (OP) describes an interpersonal situation and commenters judge the OP's conduct using tags such as YTA (``You're the Asshole'') and NTA (``Not the Asshole''). These explicit verdicts allow us to identify minority judgments relative to each post's dominant position and distinguish blaming from defending dissent. Section~\ref{subsec:method} defines the verdict grouping, alignment labels, and agreement thresholds; Appendix~\ref{sec:appendix-tags} explains the other verdict tags and variants.

After establishing the majority--minority baseline, we ask:

\begin{description}
    \item[\textbf{RQ1: Reception.}] Within high-agreement posts, how do blaming and defending dissent differ in the evaluations and conversational responses they receive?
    \item[\textbf{RQ2: Temporal dynamics.}] How does the prevalence of blaming and defending dissent change as discussions unfold in high-agreement posts?
\end{description}

% Paragraph 5: Introduce AITAH as a subsequent extension, not an institutional experiment.
We then extend the analysis to r/AITAH (AITAH), another community devoted to interpersonal moral judgment, to examine which patterns recur and which differ in another community context:

\begin{description}
    \item[\textbf{RQ3: Community context.}] Which patterns in the reception and temporal prevalence of blaming and defending dissent recur in AITAH, and which differ?
\end{description}

This comparison broadens the empirical scope without assuming that the two communities follow the same trajectory. It does not isolate the effects of any particular moderation practice or verdict rule.

% Paragraph 6: Summarize evidence and contributions without elevating baseline replication.
Our analyses combine comment scores and reply structure with annotations of direct replies. We separate OP replies from those of other commenters when examining challenge, allowing us to characterize responses beyond the person whose conduct is under judgment. The reception analyses distinguish the familiar majority--minority score--reply pattern from differences between the two directions of dissent: blaming minority judgments receive lower mean scores and greater reply engagement than defending minority judgments, while annotated replies show that minority judgments attract more challenge than majority judgments across the analyzed strata. The temporal analyses examine minority share, distinguishing cumulative prevalence from the composition of comments arriving in successive time windows. A supplementary annotation analysis compares expressed linguistic certainty across dissent directions over the full observed period, separately from RQ2. The AITAH extension examines whether reception and participation patterns recur in another community.

The central contribution is a direction-sensitive account of minority moral dissent: we distinguish challenging a crowd's defense of someone from challenging its condemnation, and examine both the reception of these positions and their prevalence as discussions unfold. The community extension identifies shared reception patterns alongside variation in participation. Together, these analyses show why understanding minority voices requires attention to what their disagreement does in a particular moral discussion, rather than treating all opposition to the majority as a single experience.
